\subsection{DashMovement a aplikace rychlosti během dashe}

Po dobu trvání dashe se ve \texttt{FixedUpdate()} přepisuje rychlost \texttt{Rigidbody} na konstantní hodnotu \texttt{dashDir * dashSpeed}. Toto zajišťuje, že hráč se pohybuje konstantní rychlostí bez ohledu na vstup nebo gravitaci. Vzorec je \texttt{rb.linearVelocity = \newline new Vector2(dashDir.x * dashSpeed, dashDir.y * dashSpeed)}.

Po uplynutí \texttt{dashTime} se dash ukončí voláním \texttt{EndDash()}. Zároveň se nastaví cooldown timer na hodnotu \texttt{dashCooldown}. Po celou dobu trvání dashe se \texttt{Rigidbody} chová jako by měl nulovou hmotnost, protože rychlost je přepisována každý snímek. To znamená, že kolize během dashe fungují správně, ale hráč je řízen čistě logikou dashe.

\begin{lstlisting}[language=CSharp,caption={Listing 10: Metoda DashMovement}]
private void FixedUpdate()
{
    if (isDashing)
    {
        DashMovement();
    }
}

private void DashMovement()
{
    if (rb == null || data == null) return;

    dashTimer -= Time.fixedDeltaTime;
    rb.linearVelocity = dashDir * data.dashSpeed;

    if (dashTimer <= 0)
    {
        isDashing = false;
    }
}

public bool IsDashing()
{
    return isDashing;
}
\end{lstlisting}